\section{Proofs}
\label{sec:appendix_proofs}

\subsection{Proof of Theorem~\ref{thm:mm1-smoothing}}
\label{section:mm1-smoothing-proof}
We focus on the case where $x_k\geq1$. Then,
\begin{align*}
x_{k+1} & =x_{k}+1\{\tau_k^A<w_k/\mu\}-1\{\tau_k^A>w_k/\mu\}\\
 & =x_{k}+\frac{e^{-\beta \tau_k^A}}{e^{-\beta \tau_k^A}+e^{-\beta w_k/\mu}}-\frac{e^{-\beta w_k/\mu}}{e^{-\beta \tau_k^A}+e^{-\beta w_k/\mu}}\\
 & =x_{k}+\frac{e^{-\beta \tau_k^A}-e^{-\beta w_k/\mu}}{e^{-\beta \tau_k^A}+e^{-\beta w_k/\mu}}.
\end{align*}
Since the inter-arrival times and workloads are exponentially distributed, by the memoryless property, we have $\tau_k^A\sim\mathsf{Exp}\left(\lambda\right)$ and
$w_k\sim \mathsf{Exp}\left(1\right)$.


The true gradient is
\[
\frac{d}{d\mu}\mathbb{E}[x_{k+1}-x_{k}]=\frac{d}{d\mu}\frac{\lambda-\mu}{\lambda+\mu}=-\frac{2\lambda}{(\lambda+\mu)^{2}}.
\]
Under our $\mathsf{softmin}_{\beta}$ approximation for the event-selection, 
%the one-step transition is
%\[
%\mathbb{E}[x_{t+1}-x_{t}]=\frac{\lambda-\mu}{\lambda+\mu}
%\]
%the $\pathwise$ gradient can be expressed as
we have
\begin{align*}
 & \mathbb{E}\left[\frac{d}{d\mu}\frac{e^{-\beta \tau_k^A}-e^{-\beta w_k/\mu}}{e^{-\beta \tau_k^A}+e^{-\beta w_k/\mu}}\right]\\
 %& =\mathbb{E}\left[\frac{d}{d\mu}\frac{e^{-\beta \tau_k^A}-e^{-\beta w_k/\mu}}{e^{-\beta \tau_k^A}+e^{-\beta w_k/\mu}}\right]\\
 & =\mathbb{E}\left[-2\beta\frac{e^{-\beta(\tau_k^A+w_k/\mu)}}{(e^{-\beta \tau_k^A}+e^{-\beta w_k/\mu})^{2}}\frac{w_k}{\mu^{2}}\right]\\
 %& =-2\frac{\beta}{\mu}\mathbb{E}\left[S_{i}\left(\frac{e^{\beta(T_{i}+S_{i})}}{(e^{\beta T_{i}}+e^{\beta S_{i}})^{2}}1\{T_{i}<S_{i}\}+\frac{e^{\beta(T_{i}+S_{i})}}{(e^{\beta T_{i}}+e^{\beta S_{i}})^{2}}1\{T_{i}>S_{i}\}\right)\right]\\
 %& =-2\frac{\beta}{\mu}\mathbb{E}\left[S_{i}\left(\frac{e^{\beta(T_{i}-S_{i})}}{(e^{\beta(T_{i}-S_{i})}+1)^{2}}1\{T_{i}<S_{i}\}+\frac{e^{\beta(S_{i}-T_{i})}}{(e^{\beta(S_{i}-T_{i})}+1)^{2}}1\{T_{i}>S_{i}\}\right)\right]\\
 & =-2\frac{\beta}{\mu}\mathbb{E}\left[\tau_k^S\left(\frac{e^{\beta(\tau_k^A-\tau_k^S)}}{(e^{\beta(\tau_k^A-\tau_k^S)}+1)^{2}}1\{\tau_k^A<\tau_k^S\}+\frac{e^{\beta(\tau_k^S-\tau_k^A)}}{(e^{\beta(\tau_k^S-\tau_k^A)}+1)^{2}}1\{\tau_k^A>\tau_k^S\}\right)\right],
\end{align*}
for $\tau_k^S=w_k/\mu$.
%
% as
% \[
% \frac{e^{\beta(T_{i}+S_{i})}}{(e^{\beta T_{i}}+e^{\beta S_{i}})^{2}}=\frac{e^{\beta(T_{i}+S_{i})}}{e^{2\beta S_{i}}(e^{\beta(T_{i}-S_{i})}+1)^{2}}=\frac{e^{\beta(T_{i}-S_{i})}}{(e^{\beta(T_{i}-S_{i})}+1)^{2}}
% \]
%
Next, note that
\begin{align*}
 & \mathbb{E}\left[\tau_k^S\frac{e^{\beta(\tau_k^A-\tau_k^S)}}{(e^{\beta(\tau_k^A-\tau_k^S)}+1)^{2}}1\{\tau_k^A<\tau_k^S\}\right]\\
 & =\mathbb{E}\left[\left.\mathbb{E}\left[\tau_k^S\frac{e^{\beta(\tau_k^A-\tau_k^S)}}{(e^{\beta(\tau_k^A-\tau_k^S)}+1)^{2}}1\{\tau_k^A<\tau_k^S\}\right]\right |\tau_k^A=t\right]\\
 & =\mathbb{E}\left[\left.\mathbb{E}\left[\tau_k^S\frac{e^{\beta(t-\tau_k^S)}}{(e^{\beta(t-\tau_k^S)}+1)^{2}}\right|\tau_k^A=t,\tau_k^S>t\right]\mathbb{P}(\tau_k^S>t|\tau_k^A=t)\right]\\
 %& =\mathbb{E}\left[\mathbb{E}\left[(t+S')\frac{e^{-\beta S'}}{(e^{-\beta S'}+1)^{2}}\right]\mathbb{P}(\tau_k^S>t|\tau_k^A=t)\right]\\
 & =\mathbb{E}\left[\mathbb{E}\left[(t+S')\frac{e^{-\beta S'}}{(e^{-\beta S'}+1)^{2}}\right]\mathbb{P}(\tau_k^S>t|\tau_k^A=t)\right] \mbox{ for $S'\sim \mathsf{Exp}(\mu)$}\\
 & =\mathbb{E}\left[(tA(\beta,\mu)+B(\beta,\mu))\mathbb{P}(\tau_k^S>t|\tau_k^A=t)\right]\\
 & =\mathbb{E}\left[(\tau_k^A A(\beta,\mu)+B(\beta,\mu))e^{-\mu \tau_k^A}\right]\\
 & =\frac{\lambda}{(\lambda+\mu)^{2}}A(\beta,\mu)+\frac{\lambda}{(\lambda+\mu)}B(\beta,\mu)
\end{align*}
where
\begin{align*}
A(\beta,\mu) & =\frac{\mu\left(\beta-\mu H\left(\frac{\mu}{2\beta}\right)+\mu H\left(\frac{\mu}{2\beta}-\frac{1}{2}\right)\right)}{2\beta^{2}}\\
 & =\frac{\mu}{2\beta}+\frac{\mu^{2}}{2\beta^{2}}\left(\underbrace{H\left(\frac{\mu}{2\beta}-\frac{1}{2}\right)-H\left(\frac{\mu}{2\beta}\right)}_{\tilde{H}(\beta,\mu)}\right)\\
B(\beta,\mu) & =\frac{\mu}{4\beta^{3}}\left(2\beta H\left(\frac{\mu}{2\beta}\right)-2\beta H\left(\frac{\mu}{2\beta}-\frac{1}{2}\right)-\mu\psi^{(1)}\left(\frac{\beta+\mu}{2\beta}\right)+\mu\psi^{(1)}\left(\frac{2\beta+\mu}{2\beta}\right)\right)\\
 & =-\frac{\mu}{2\beta^{2}}\tilde{H}(\beta,\mu)+\frac{\mu^{2}}{4\beta^{3}}\left(\underbrace{\psi^{(1)}\left(\frac{2\beta+\mu}{2\beta}\right)-\psi^{(1)}\left(\frac{\beta+\mu}{2\beta}\right)}_{\tilde{\psi}^{(1)}(\beta,\mu)}\right)
\end{align*}
Moreover, note that
\begin{align*}
H\left(\frac{\mu}{2\beta}\right) & =\psi^{(0)}\left(\frac{\mu}{2\beta}+1\right)+\gamma\\
H\left(\frac{\mu}{2\beta}-\frac{1}{2}\right) & =\psi^{(0)}\left(\frac{\mu}{2\beta}+\frac{1}{2}\right)+\gamma
\end{align*}
Note that
\begin{align*}
H\left(\frac{\mu}{2\beta}\right)-H\left(\frac{\mu}{2\beta}-\frac{1}{2}\right) & =\log\left(\frac{\mu}{2\beta}+1\right)+\frac{1}{\frac{\mu}{2\beta}+1}\\
 & -\log\left(\frac{\mu}{2\beta}+\frac{1}{2}\right)+\frac{1}{\frac{\mu}{2\beta}+1}.
\end{align*}

Similarly,
\begin{align*}
 & \mathbb{E}\left[\tau_k^S\frac{e^{\beta(\tau_k^S-\tau_k^A)}}{(e^{\beta(\tau_k^S-\tau_k^A)}+1)^{2}}1\{\tau_k^A>\tau_k^S\}\right]\\
 & =\mathbb{E}\left[\mathbb{E}\left[\left.s\frac{e^{-\beta(\tau_k^A-s)}}{(e^{-\beta(\tau_k^A-s)}+1)^{2}}1\{\tau_k^A>s\}\right|\tau_k^S=s\right]\right]\\
 & =\mathbb{E}\left[\mathbb{E}\left[\left.s\frac{e^{-\beta(\tau_k^A-s)}}{(e^{-\beta(\tau_k^A-s)}+1)^{2}}\right|\tau_k^A>s,\tau_k^S=s\right]\mathbb{P}(\tau_k^A>s|\tau_k^S=s)\right]\\
 & =\mathbb{E}\left[\mathbb{E}\left[s\frac{e^{-\beta T'}}{(e^{-\beta T'}+1)^{2}}\right]\mathbb{P}(\tau_k^A>s|\tau_k^S=s)\right] \mbox{ for $T'\sim \mathsf{Exp}(\lambda)$}\\
 & =\mathbb{E}\left[\tau_k^S A(\beta,\lambda)e^{-\lambda \tau_k^S}\right]\\
 & =\frac{\mu}{(\lambda+\mu)^{2}} A(\beta,\lambda).
\end{align*}
Then,
\begin{align*}
% & -2\frac{\beta}{\mu}\mathbb{E}\left[\tau_k^S\left(e^{\beta(\tau_k^A-\tau_k^S)}1\{\tau_k^A<\tau_k^S\}+e^{\beta(\tau_k^S-\tau_k^A)}1\{\tau_k^A>\tau_k^S\}\right)\right]\\
 &-2\frac{\beta}{\mu}\mathbb{E}\left[\tau_k^S\left(\frac{e^{\beta(\tau_k^A-\tau_k^S)}}{(e^{\beta(\tau_k^A-\tau_k^S)}+1)^{2}}1\{\tau_k^A<\tau_k^S\}+\frac{e^{\beta(\tau_k^S-\tau_k^A)}}{(e^{\beta(\tau_k^S-\tau_k^A)}+1)^{2}}1\{\tau_k^A>\tau_k^S\}\right)\right]\\
 & =-2\frac{\beta}{\mu}\left(\frac{\lambda}{(\lambda+\mu)^{2}}A(\beta,\mu)+\frac{\lambda}{(\lambda+\mu)}B(\beta,\mu)+\frac{\mu}{(\lambda+\mu)^{2}}A(\beta,\lambda)\right)\\
 & =-2\frac{\beta}{\mu}\left(\frac{\lambda A(\beta,\mu)+\mu A(\beta,\lambda)}{(\lambda+\mu)^{2}}+\frac{\lambda}{(\lambda+\mu)}B(\beta,\mu)\right)\\
 & =\frac{-2\lambda}{(\lambda+\mu)^{2}}-\frac{2\beta}{\mu}\left(\frac{\lambda\mu^{2}\tilde{H}(\beta,\mu)+\mu\lambda^{2}\tilde{H}(\beta,\lambda)}{2\beta^{2}(\lambda+\mu)^{2}}-\frac{\lambda\mu}{2\beta^{2}(\lambda+\mu)}\tilde{H}(\beta,\mu)+\frac{\lambda}{(\lambda+\mu)}\frac{\mu^{2}}{4\beta^{3}}\tilde{\psi}^{(1)}(\beta,\mu)\right).
\end{align*}
Note that as $\beta\to\infty$,
\begin{align*}
\lim_{\beta\to\infty}\tilde{H}(\beta,\mu) & =\gamma+\psi^{(0)}\left(\frac{1}{2}\right),\\
\lim_{\beta\to\infty}\tilde{\psi}^{(1)}(\beta,\mu) & =-\frac{\pi^{2}}{3}.
\end{align*}
This means that the leading order term is $O(1/\beta^{2})$. In particular,
\begin{align*}
-\frac{2\beta}{\mu}\left(\frac{\lambda\mu^{2}\tilde{H}(\beta,\mu)+\mu\lambda^{2}\tilde{H}(\beta,\lambda)}{2\beta^{2}(\lambda+\mu)^{2}}-\frac{\lambda\mu\tilde{H}(\beta,\mu)}{2\beta^{2}(\lambda+\mu)}\right) & \sim\frac{\pi^{2}\lambda^{2}(\mu-\lambda)}{6\beta^{2}(\lambda+\mu)^{2}}
\end{align*}
Finally, we have the second-order term
\[
-\frac{2\beta}{\mu}\frac{\lambda}{(\lambda+\mu)}\frac{\mu^{2}}{4\beta^{3}}\tilde{\psi}^{(1)}(\beta,\mu)\sim\frac{\pi^{2}\lambda\mu}{6\beta^{2}(\lambda+\mu)}.
\]
Thus, we have the following characterization of the bias:
\[
\mathbb{E}\left[\frac{d}{d\mu}\frac{e^{-\beta \tau_k^A}-e^{-\beta w_k/\mu}}{e^{-\beta \tau_k^A}+e^{-\beta w_k/\mu}}\right]-\left(\frac{-2\lambda}{(\lambda+\mu)^{2}}\right)\sim\frac{1}{\beta^{2}}\frac{\pi^{2}\lambda(\mu^{2}-\lambda^{2}+2\mu\lambda)}{6(\lambda+\mu)^{2}}+o\left(\frac{1}{\beta^{2}}\right)
\]

For variance, we have
\begin{align*}
 & \mathbb{E}\left[\left(-2\beta\frac{e^{\beta(\tau_k^A+w_k/\mu)}}{(e^{\beta \tau_k^A}+e^{\beta w_k/\mu})^{2}}\frac{w_k}{\mu^{2}}\right)^{2}\right]\\
 & =\mathbb{E}\left[\frac{4\beta^{2}}{\mu^{2}}\frac{e^{2\beta(\tau_k^A+\tau_k^S)}}{\left(e^{\beta \tau_k^A}+e^{\beta \tau_k^S}\right)^{4}}\tau_k^{S,2}\right]\\
 & =\frac{4\beta^{2}}{\mu^{2}}\mathbb{E}\left[\tau_k^{S,2}\left(\frac{e^{2\beta(\tau_k^A+\tau_k^S)}}{\left(e^{\beta \tau_k^A}+e^{\beta \tau_k^S}\right)^{4}}1\{\tau_k^A<\tau_k^S\}+\frac{e^{2\beta(\tau_k^A+\tau_k^S)}}{\left(e^{\beta \tau_k^A}+e^{\beta \tau_k^S}\right)^{4}}1\{\tau_k^A>\tau_k^S\}\right)\right]\\
 & =\frac{4\beta^{2}}{\mu^{2}}\mathbb{E}\left[\tau_k^{S,2}\left(\frac{e^{2\beta(\tau_k^A-\tau_k^S)}}{\left(e^{\beta(\tau_k^A-\tau_k^S)}+1\right)^{4}}1\{\tau_k^A<\tau_k^S\}+\frac{e^{2\beta(\tau_k^S-\tau_k^A)}}{\left(e^{\beta(\tau_k^S-\tau_k^A)}+1\right)^{4}}1\{\tau_k^A>\tau_k^S\}\right)\right].
\end{align*}
Note that
\begin{align*}
 & \mathbb{E}\left[\tau_k^{S,2}\frac{e^{2\beta(\tau_k^A-\tau_k^S)}}{\left(e^{\beta(\tau_k^A-\tau_k^S)}+1\right)^{4}}1\{\tau_k^A<\tau_k^S\}\right]\\
 & =\mathbb{E}\left[\mathbb{E}\left[\left.\tau_k^{S,2}\frac{e^{2\beta(\tau_k^A-\tau_k^S)}}{\left(e^{\beta(\tau_k^A-\tau_k^S)}+1\right)^{4}}1\{t<\tau_k^S\}\right]\right|\tau_k^A=t\right]\\
 %& =\mathbb{E}\left[\mathbb{E}\left[S_{i}^{2}\frac{e^{2\beta(t-S_{i})}}{(e^{\beta(t-S_{i})}+1)^{2}}|T_{i}=t,S_{i}>t\right]\mathbb{P}(S_{i}>t|T_{i}=t)\right]\\
 & =\mathbb{E}\left[\mathbb{E}\left[(t+S')^{2}\frac{e^{-2\beta S'}}{(e^{-\beta S'}+1)^{4}}\right]\mathbb{P}(\tau_k^S>t|\tau_k^A=t)\right] \mbox{ for $S'\sim \mathsf{Exp}(\mu)$}\\
 & =\mathbb{E}\left[\mathbb{E}\left[(t^{2}+2S'+S'^{2})\frac{e^{-2\beta S'}}{(e^{-\beta S'}+1)^{4}}\right]\mathbb{P}(S_{i}>t|T_{i}=t)\right]\\
 & =\mathbb{E}\left[\left(t^{2}\tilde{A}(\beta,\mu)+t\tilde{B}(\beta,\mu)+\tilde{C}(\beta,\mu)\right)\mathbb{P}(S_{i}>t|T_{i}=t)\right]\\
 & =\mathbb{E}\left[\left(T_{i}^{2}\tilde{A}(\beta,\mu)+T_{i}\tilde{B}(\beta,\mu)+\tilde{C}(\beta,\mu)\right)e^{-\mu T_{i}}\right]\\
 & =\frac{2\lambda}{(\lambda+\mu)^{3}}\tilde{A}(\beta,\mu)+\frac{\lambda}{(\lambda+\mu)^{2}}\tilde{B}(\beta,\mu)+\frac{\lambda}{(\lambda+\mu)}\tilde{C}(\beta,\mu).
\end{align*}
Similarly,
\begin{align*}
 & \mathbb{E}\left[\tau_k^{S,2}\frac{e^{2\beta(\tau_k^S-\tau_k^A)}}{\left(e^{\beta(\tau_k^S-\tau_k^A)}+1\right)^{4}}1\{\tau_k^A>\tau_k^S\}\right]\\
 & =\mathbb{E}\left[\mathbb{E}\left[\left.s^{2}\frac{e^{2\beta(s-\tau_k^A)}}{\left(e^{\beta(s-\tau_k^A)}+1\right)^{4}}1\{\tau_k^A>s\}\right|\tau_k^S=s\right]\right]\\
 & =\mathbb{E}\left[\mathbb{E}\left[s^{2}\frac{e^{-2\beta(\tau_k^A-s)}}{(e^{-\beta(\tau_k^A-s)}+1)^{4}}|\tau_k^A>s,\tau_k^S=s\right]\mathbb{P}(\tau_k^A>s|\tau_k^S=s)\right]\\
 & =\mathbb{E}\left[\mathbb{E}\left[s^{2}\frac{e^{-2\beta T'}}{(e^{-\beta T'}+1)^{4}}\right]\mathbb{P}(\tau_k^A>s|\tau_k^S=s)\right] \mbox{ for $T'\sim \mathsf{Exp}(\lambda)$}\\
 & =\mathbb{E}\left[\tilde{A}(\beta,\lambda)\tau_k^{S,2}e^{-\lambda \tau_k^S}\right]\\
 & =\frac{2\mu}{(\lambda+\mu)^{2}}\tilde{A}(\beta,\lambda).
\end{align*}
Putting the above two parts together, we have
\begin{align*}
 %& \frac{4\beta^{2}}{\mu^{2}}\mathbb{E}\left[S_{i}^{2}\left(\frac{e^{\beta(T_{i}-S_{i})2}}{(e^{\beta(T_{i}-S_{i})}+1)^{4}}1\{T_{i}<S_{i}\}+\frac{e^{\beta(S_{i}-T_{i})2}}{(e^{\beta(S_{i}-T_{i})}+1)^{4}}1\{T_{i}>S_{i}\}\right)\right]\\
 &\frac{4\beta^{2}}{\mu^{2}}\mathbb{E}\left[\tau_k^{S,2}\left(\frac{e^{2\beta(\tau_k^A-\tau_k^S)}}{\left(e^{\beta(\tau_k^A-\tau_k^S)}+1\right)^{4}}1\{\tau_k^A<\tau_k^S\}+\frac{e^{2\beta(\tau_k^S-\tau_k^A)}}{\left(e^{\beta(\tau_k^S-\tau_k^A)}+1\right)^{4}}1\{\tau_k^A>\tau_k^S\}\right)\right]\\
 & =\frac{4\beta^{2}}{\mu^{2}}\left(\frac{2\lambda}{(\lambda+\mu)^{3}}\tilde{A}(\beta,\mu)+\frac{\lambda}{(\lambda+\mu)^{2}}\tilde{B}(\beta,\mu)+\frac{\lambda}{(\lambda+\mu)}\tilde{C}(\beta,\mu)+\frac{2\mu}{(\lambda+\mu)^{2}}\tilde{A}(\beta,\lambda)\right)\\
 & \sim\frac{4\beta\lambda}{3\mu(\lambda+\mu)^{2}} \mbox{ as $\beta\to\infty$.}
\end{align*}


\subsection{Proof of Theorem~\ref{thm:mm1_variance}}
\label{sec:proof_mm1_variance}

By assumption, $h \leq c(\mu - \lambda)$ for some $c < 1$. First, we develop bound for $\var_{\infty} (\hat{\nabla}^{\mathsf{R}} J_{N}(\theta; \xi_{1:N}))$.
%= \Theta\left(N^{-1} \theta^{-2} (1-\rho)^{-4} \right)$. 
We can compute the variance by conditioning on the value of $Y$:
\begin{align*}
    \var_{\infty}(\hat{Q}_{N}(\mu-hY)\cdot(\log Y-\frac{1}{\theta}))&=\mathbb{E}\left[\var_{\infty}\left(\hat{Q}_{N}(\mu-hy)(\log y-\frac{1}
{\theta})|Y=y\right)\right]\\
&+\var\left(\mathbb{E}\left[\hat{Q}_{N}(\mu-hy)(\log y-\frac{1}{\theta})|Y=y)\right]\right).
\end{align*}
For the first term, note that the asymptotic variance in the CLT for the ergodic estimator $\hat{Q}_{N}(\mu)$ is $\var_{\infty}(\hat Q(\mu))=\frac{2\rho(1+\rho)}{(1-\rho)^{4}}$. Then, we have $\var_{\infty}(\hat{Q}_{N}(\mu)) = \frac{2\rho(1+\rho)}{N(1-\rho)^{4}}$. Since $h \leq \mu$, $\var_{\infty}(\hat{Q}_{N}(\mu))= \frac{2\rho(1+\rho)}{N(1-\rho)^{4}}$.
Then,
\begin{align*}
&\E \left[ \var_{\infty}\left(\hat{Q}_{N}(\mu-hy)\cdot(\log y+\frac{1}{\theta})|Y = y\right) \right] \\
& =\E \left[  (\log y+\frac{1}{\theta})^{2} \var_{\infty}(\hat{Q}_{N}(\mu-hy)|Y = y) \right]\\
& \geq \E \left[  (\log Y+\frac{1}{\theta})^{2}
\frac{2\rho(1+\rho)}{N(1-\rho)^{4}} \right]\\
& = \frac{1}{\theta^{2}}
\frac{2\rho(1+\rho)}{N(1-\rho)^{4}}=\Theta\left( (1-\rho)^{-4} \right),
\end{align*}
where the last equality uses the fact that for $Y\sim\text{Beta}(\theta,1)$,
\[
\E[\log Y]=\psi(\theta) - \psi(\theta + 1)=-\frac{1}{\theta}
\]
and
\[
\E \left[  (\log Y+\frac{1}{\theta})^{2}\right]=\var(\log Y) = \psi_{1}(\theta) - \psi_{1}(\theta + 1) = \frac{1}{\theta^{2}}.
\]
For the second term, we plug in the true estimand as the expectation of $Q_{N}(\mu)$, i.e. $Q(\mu) = \frac{\lambda}{\mu - \lambda}$. Then,
\begin{align*}
&\var\left(\mathbb{E}\left[\hat{Q}_{N}(\mu-hy)(\log y+\frac{1}{\theta})|Y=y)\right]\right) \\
& = \var\left(\frac{\lambda}{\mu - hY - \lambda}
\left( \log Y + \frac{1}{\theta} \right)\right) \\
& = \E \left[ \left(\frac{\lambda}{\mu - hY - \lambda} \right)^{2}
\left( \log Y + \frac{1}{\theta} \right)^{2} \right]
- \E \left[ \left(\frac{\lambda}{\mu - hY - \lambda} \right)
\left( \log Y + \frac{1}{\theta} \right) \right]^{2}.
\end{align*}
We proceed to evaluate these expectations analytically.
\begin{align*}
& \E \left[ \left(\frac{\lambda}{\mu - hY - \lambda} \right)
\left( \log Y + \frac{1}{\theta} \right) \right] \\
&= \frac{\rho}{1-\rho}
\left(
\frac{\Gamma(\theta)}{\Gamma(1+\theta)}
F^{2}_{1}
\left(1,\theta,1+\theta,\frac{h}{\mu-\lambda} \right)
-
\theta\Phi \left(\frac{h}{\mu - \lambda},2,\theta \right)
\right) \\
&=O((1-\rho)^{-1}),
\end{align*}
where $F^{2}_{1}$ is the Hypergoemetric 2F1  function and $\Phi$ is the Lerch transcendental function.
We also have
\begin{align*}
& \E \left[ \left(\frac{\lambda}{\mu - hY - \lambda} \right)^{2}
\left( \log Y + \frac{1}{\theta} \right)^{2} \right] \\
& = \frac{\lambda^{2}}{\theta^{2}(\mu - \lambda)^{3}}
\left(
2(\mu - \lambda)
+ h\theta^{3}\Phi \left(\frac{h}{\mu - \lambda},2,\theta+1 \right)
+ h\theta^{3}(\theta - 1)\Phi \left(\frac{h}{\mu - \lambda},3,\theta+1 \right) \right. \\
& 
+ (\mu - \lambda) 
\left( \theta\frac{\mu - \lambda}{\mu - h - \lambda} - (\theta - 1) F_{1}^{2}(1,\theta,1+\theta,\frac{h}{\mu-\lambda})
\right) \\
&\left. +2(\mu -\lambda) F_{2}^{3}\left((2,\theta,\theta),(1+\theta,1+\theta),\frac{h}{\mu-\lambda}\right)
\right) \\
& = O((1-\rho)^{-2}).
\end{align*}
Taking the difference between the above two parts under the limit as $(1-\rho)\to 0$, we have
\[
\var\left(\frac{\lambda}{\mu - hY - \lambda}
\left( \log Y + \frac{1}{\theta} \right)\right) =O((1-\rho)^{-2}).
\]

Next, we develop a bound for $\var_{\infty} (\hat{\nabla} J_{N}(\theta; \xi_{1:N}))$.
\begin{align*}
    \var_{\infty}\left(
    h \cdot \hat{\nabla} Q_{N}(\mu - h Y) \left( \frac{1}{\theta} Y \log Y \right)
    \right)&=
    \E\left[\var_{\infty}\left(h \cdot \hat{\nabla} Q_{N}(\mu - h Y) \left( \frac{1}{\theta} Y \log Y \right)|Y=y\right)\right]\\
&+\var\left(\mathbb{E}\left[h \cdot \hat{\nabla} Q_{N}(\mu - h Y) \left( \frac{1}{\theta} Y \log Y \right)|Y = y\right]\right)
\end{align*}
For the first term, we can use the fact that $|Y\log Y| \leq 1/e$ almost surely since $Y\in[0,1]$. We also use recent results in~\cite{hu2023comparison}, which compute the asymptotic variance of the IPA estimator:
\[
\var_{\infty}(\hat{\nabla} Q_{N}(\mu)) = \frac{1 + 16\rho + 27 \rho^{2} + 2 \rho^{3} + 6\rho^{4}}{\mu^{2}N(1 + \rho)(1-\rho)^{5}} \leq 52\mu^{-2}N^{-1}(1-\rho)^{-5}
\]
Under $\mu - hy$, the congestion factor $1 - \frac{\lambda}{\mu - hy} = \frac{\mu - hy - \lambda}{\mu - hy} \geq \frac{\mu - cy(\mu - \lambda) - \lambda}{\mu} = (1-cy)(1-\rho)$ and $\mu-hy\geq \mu-h\geq (1-c)\mu$.
So we have the bound,
\begin{align*}
 &\var_{\infty}\left(h\hat{\nabla} Q_{N}(\mu - h Y) \left( \frac{1}{\theta} y \log y\right)|Y=y\right) \\
 & \leq  h^{2} \theta^{-2}(y \log y)^{2}\var_{\infty}(\hat{\nabla} Q_{N}(\mu - h y)) \\
 & \leq h^{2}\mu^{-2}\theta^{-2}e^{-2} 52N^{-1}(1-c)^{-7}(1-\rho)^{-5} \\
 & = O(N^{-1} h^{2} \mu^{-2} (1-\rho)^{-5}) \\
 &= O(N^{-1}  (1-\rho)^{-3})
\end{align*}
since $h = O(1-\rho)$.

For the second term, we plug in the true estimand as the mean of $\hat{\nabla} Q_{N}(\mu)$, i.e., $ \nabla Q(\mu) = -\frac{\rho}{\mu(1-\rho)^{2}}$,
\begin{align*}
&\E\left[h \cdot \hat{\nabla} Q_{N}(\mu - h Y) \left( \frac{1}{\theta} Y \log Y \right)|Y = y\right] 
= - h\frac{1}{\theta} (y \log y) \frac{\lambda }{(\mu - hy - \lambda)^{2}}.
\end{align*}
% Note that for any $y\in [0,1]$ this has lower bound $0$ and upper bound $h\theta^{-1} e^{-1} \frac{\rho}{(\mu - h)(1-c)^{2}(1-\rho)^{2}}$, which gives a variance upper bound of
We next evaluate the variance analytically,
\begin{align*}
&\var\left(
- h\frac{1}{\theta} (Y \log Y) \frac{\lambda }{(\mu - hY - \lambda)^{2}}
\right) \\
&= 
\E\left[ 
\left(- h\frac{1}{\theta} (Y \log Y) \frac{\lambda }{(\mu - hY - \lambda)^{2}}
\right)^{2}
\right] 
- \E\left[ 
- h\frac{1}{\theta} (Y \log Y) \frac{\lambda }{(\mu - hY - \lambda)^{2}}
\right]^{2}
\end{align*}
Since
\begin{align*}
&\E\left[ 
- h\frac{1}{\theta} (Y \log Y) \frac{\lambda }{(\mu - hY - \lambda)^{2}}
\right] \\
&= \frac{h}{(1+\theta)^{2}}
\frac{\lambda}{(\mu - \lambda)^{2}}
F^{3}_{2} 
\left((2,1+\theta,1+\theta), 
(2+\theta, 2+ \theta), \frac{h}{\mu -\lambda} \right) \\
& = O((1-\rho)^{-1})
\end{align*}
and
\begin{align*}
&\E\left[ 
\left(- h\frac{1}{\theta} (Y \log Y) \frac{\lambda }{(\mu - hY - \lambda)^{2}}
\right)^{2}
\right] \\
&=
2h\theta^{2}\Gamma(\theta)^{3}\Gamma(2+\theta)^{-3}\frac{\lambda^{2}}{(\mu - \lambda)^{2}} \times 
\left( F^{4}_{3} 
\left((3,1+\theta,1+\theta,1+\theta), 
(2+\theta, 2+ \theta,2+ \theta), \frac{h}{\mu -\lambda} \right) \right. \\
& \left. -F^{4}_{3} 
\left((4,1+\theta,1+\theta,1+\theta), 
(2+\theta, 2+ \theta,2+ \theta), \frac{h}{\mu -\lambda} \right) \right) \\
& = O((1-\rho)^{-2}),
\end{align*}
we have
\[
\var\left(
- h\frac{1}{\theta} (Y \log Y) \frac{\lambda }{(\mu - hY - \lambda)^{2}}
\right) = O( (1-\rho)^{-2}).
\]


\subsection{Proof of Corollary~\ref{cor:reinforce-baseline}}
\label{sec:proof_reinforce_baseline}

First, we can explicitly characterize the optimal baseline:
\begin{align*}
b^{*} &= \frac{\E[Q(\mu -hY) \nabla_{\theta} \log \pi_{\theta}(Y)^{2}]}{\E[\nabla_{\theta}\log \pi_{\theta}(Y)^{2}]} \\
 &= \frac{\E[Q(\mu -hY) \left( \log Y + 1/\theta \right)^{2}]}{\E[\left( \log Y + 1/\theta \right)^{2}]} \\
&= \frac{\lambda}{\mu - \lambda}\underbrace{\left[F^{2}_{1}\left(1,\theta,1+\theta,\frac{h}{\mu-\lambda}\right) -2\theta^{2}\Phi\left(\frac{h}{\mu-\lambda},2,\theta\right) + 2\frac{h}{\mu-\lambda}\theta^{3}\Phi\left(\frac{h}{\mu-\lambda},3,1+\theta\right)\right]}_{b(\theta)}  \\
& = O((1-\rho)^{-1}).
\end{align*}
Next, we plug this into the $\reinforce$ estimator. 
\begin{align*}
    \var_{\infty}((\hat{Q}_{N}(\mu-hY)-b^*)\cdot(\log Y-\frac{1}{\theta}))&=\mathbb{E}\left[\var_{\infty}\left((\hat{Q}_{N}(\mu-hy)-b^*)(\log y-\frac{1}
{\theta})|Y=y\right)\right]\\
&+\var\left(\mathbb{E}\left[(\hat{Q}_{N}(\mu-hy) - b^{*})(\log y-\frac{1}{\theta})|Y=y\right]\right).
\end{align*}
Note that since $b^{*}$ is a constant, the first term, i.e., the mean of the conditional variance given $Y$, has the same value as in Theorem \ref{thm:mm1_variance}. For the second term, note that since $h = c(\mu - \lambda)$,
% For the second, i.e., we can explicitly compute the variance of the conditional mean given $Y$ with Mathematica. However, this leads to an unwieldy expressio
\begin{align*}
& \var \left[ \left(\frac{\lambda}{\mu - hY - \lambda}  - b^{*}\right)
\left( \log Y + \frac{1}{\theta} \right) \right]  \\
&= \left(\frac{\lambda}{\mu - \lambda}\right)^{2} \var \left[ \left(\frac{1}{1- cY}  - b(\theta)\right)
\left( \log Y + \frac{1}{\theta} \right) \right].
\end{align*}
Since $\var \left[ \left(\frac{1}{1- cY}  - b(\theta)\right)
\left( \log Y + \frac{1}{\theta} \right) \right] > 0$ and doesn't depend on $\mu$ or $\lambda$, this confirms that the second term is $\Theta((1-\rho)^{-2})$

